\documentclass[10pt,twocolumn]{article}
\usepackage[margin=0.8in]{geometry}
\usepackage{amsmath,amssymb,booktabs,graphicx,hyperref,xcolor,framed}
\hypersetup{colorlinks=true,linkcolor=blue!50!black,urlcolor=blue!50!black}
\setlength{\columnsep}{0.28in}
\newcommand{\ActA}{18}
\newcommand{\NptsA}{240}
\newcommand{\GaussA}{1.2\times10^{-8}}
\newcommand{\GaussHiA}{9.6\times10^{-8}}
\newcommand{\BianA}{7.6\times10^{-13}}
\newcommand{\BianHiA}{1.5\times10^{-10}}
\newcommand{\AsymA}{2.9\times10^{-10}}
\newcommand{\SplitA}{1.2\times10^{-15}}
\newcommand{\RconnVsGaussA}{4.8\times10^{-10}}
\newcommand{\EigMinA}{3.4\times10^{-4}}
\newcommand{\EigMaxA}{1.31}
\newcommand{\CondA}{4849}
\newcommand{\LegacyA}{1.74\times10^{3}}
\newcommand{\RmedA}{24.1}
\newcommand{\RloA}{-1027}
\newcommand{\RhiA}{6666}
\newcommand{\AbsRA}{307}
\newcommand{\KmedA}{3.8\times10^{-5}}
\newcommand{\KloA}{-0.007}
\newcommand{\KhiA}{0.022}
\newcommand{\KNegA}{45}
\newcommand{\GeoMsA}{0.60}
\newcommand{\GeoMeA}{1.90}
\newcommand{\GeoLenA}{6.263}
\newcommand{\LineLenA}{6.623}
\newcommand{\GeoNormalA}{94}
\newcommand{\LineNormalA}{3}
\newcommand{\GeoHawkA}{0.108}
\newcommand{\LineHawkA}{0.049}
\newcommand{\GeoEinA}{0.227}
\newcommand{\LineEinA}{0.101}
\newcommand{\GeoSpeedA}{2\times10^{-4}}
\newcommand{\RDustA}{-138.4}
\newcommand{\ShCoDustA}{$<$1}
\newcommand{\ShGrDustA}{14}
\newcommand{\ShQuDustA}{54}
\newcommand{\ShHoDustA}{19}
\newcommand{\ShPlDustA}{13}
\newcommand{\RRadiationA}{43.3}
\newcommand{\ShCoRadiationA}{$<$1}
\newcommand{\ShGrRadiationA}{12}
\newcommand{\ShQuRadiationA}{42}
\newcommand{\ShHoRadiationA}{28}
\newcommand{\ShPlRadiationA}{18}
\newcommand{\RVacuumA}{501.8}
\newcommand{\ShCoVacuumA}{$<$1}
\newcommand{\ShGrVacuumA}{13}
\newcommand{\ShQuVacuumA}{50}
\newcommand{\ShHoVacuumA}{18}
\newcommand{\ShPlVacuumA}{20}
\newcommand{\RBlackholeA}{-15.9}
\newcommand{\ShCoBlackholeA}{$<$1}
\newcommand{\ShGrBlackholeA}{27}
\newcommand{\ShQuBlackholeA}{21}
\newcommand{\ShHoBlackholeA}{32}
\newcommand{\ShPlBlackholeA}{19}
\newcommand{\RWormholeA}{363.4}
\newcommand{\ShCoWormholeA}{$<$1}
\newcommand{\ShGrWormholeA}{16}
\newcommand{\ShQuWormholeA}{37}
\newcommand{\ShHoWormholeA}{26}
\newcommand{\ShPlWormholeA}{21}
\newcommand{\RHawkingA}{-8.6}
\newcommand{\ShCoHawkingA}{$<$1}
\newcommand{\ShGrHawkingA}{25}
\newcommand{\ShQuHawkingA}{29}
\newcommand{\ShHoHawkingA}{28}
\newcommand{\ShPlHawkingA}{17}
\newcommand{\BerryA}{8\times10^{-16}}
\newcommand{\ActB}{18}
\newcommand{\NptsB}{240}
\newcommand{\GaussB}{1.6\times10^{-8}}
\newcommand{\GaussHiB}{1.7\times10^{-7}}
\newcommand{\BianB}{1.1\times10^{-12}}
\newcommand{\BianHiB}{8.2\times10^{-11}}
\newcommand{\AsymB}{4.2\times10^{-10}}
\newcommand{\SplitB}{1.2\times10^{-15}}
\newcommand{\RconnVsGaussB}{4.3\times10^{-10}}
\newcommand{\EigMinB}{2.2\times10^{-4}}
\newcommand{\EigMaxB}{1.29}
\newcommand{\CondB}{7562}
\newcommand{\LegacyB}{2.62\times10^{3}}
\newcommand{\RmedB}{-27.9}
\newcommand{\RloB}{-5496}
\newcommand{\RhiB}{3773}
\newcommand{\AbsRB}{524}
\newcommand{\KmedB}{7.8\times10^{-5}}
\newcommand{\KloB}{-0.006}
\newcommand{\KhiB}{0.025}
\newcommand{\KNegB}{41}
\newcommand{\GeoMsB}{0.60}
\newcommand{\GeoMeB}{1.90}
\newcommand{\GeoLenB}{6.205}
\newcommand{\LineLenB}{6.560}
\newcommand{\GeoNormalB}{97}
\newcommand{\LineNormalB}{5}
\newcommand{\GeoHawkB}{0.118}
\newcommand{\LineHawkB}{0.041}
\newcommand{\GeoEinB}{0.415}
\newcommand{\LineEinB}{0.084}
\newcommand{\GeoSpeedB}{6\times10^{-5}}
\newcommand{\RDustB}{-384.8}
\newcommand{\ShCoDustB}{$<$1}
\newcommand{\ShGrDustB}{14}
\newcommand{\ShQuDustB}{42}
\newcommand{\ShHoDustB}{26}
\newcommand{\ShPlDustB}{18}
\newcommand{\RRadiationB}{-143.4}
\newcommand{\ShCoRadiationB}{$<$1}
\newcommand{\ShGrRadiationB}{13}
\newcommand{\ShQuRadiationB}{38}
\newcommand{\ShHoRadiationB}{30}
\newcommand{\ShPlRadiationB}{18}
\newcommand{\RVacuumB}{-304.5}
\newcommand{\ShCoVacuumB}{$<$1}
\newcommand{\ShGrVacuumB}{8}
\newcommand{\ShQuVacuumB}{37}
\newcommand{\ShHoVacuumB}{31}
\newcommand{\ShPlVacuumB}{24}
\newcommand{\RBlackholeB}{-24.6}
\newcommand{\ShCoBlackholeB}{$<$1}
\newcommand{\ShGrBlackholeB}{22}
\newcommand{\ShQuBlackholeB}{21}
\newcommand{\ShHoBlackholeB}{34}
\newcommand{\ShPlBlackholeB}{22}
\newcommand{\RWormholeB}{519.8}
\newcommand{\ShCoWormholeB}{$<$1}
\newcommand{\ShGrWormholeB}{12}
\newcommand{\ShQuWormholeB}{35}
\newcommand{\ShHoWormholeB}{23}
\newcommand{\ShPlWormholeB}{29}
\newcommand{\RHawkingB}{-3.2}
\newcommand{\ShCoHawkingB}{$<$1}
\newcommand{\ShGrHawkingB}{28}
\newcommand{\ShQuHawkingB}{17}
\newcommand{\ShHoHawkingB}{34}
\newcommand{\ShPlHawkingB}{21}
\newcommand{\BerryB}{7\times10^{-16}}
\newcommand{\UBoundA}{99.6}
\newcommand{\UStatesA}{240}
\newcommand{\UHeldA}{239}
\newcommand{\UBetaA}{0.0000}
\newcommand{\UBetaLoA}{-0.0011}
\newcommand{\UBetaHiA}{0.0010}
\newcommand{\UCorrA}{-0.029}
\newcommand{\UBoundB}{100.0}
\newcommand{\UStatesB}{240}
\newcommand{\UHeldB}{240}
\newcommand{\UBetaB}{-0.0000}
\newcommand{\UBetaLoB}{-0.0013}
\newcommand{\UBetaHiB}{0.0006}
\newcommand{\UCorrB}{-0.025}


\title{\vspace{-1.2em}\textbf{A Student-t($\nu=0.8$) Prior on Measurement Processes Yields a Riemannian\\
Manifold Holding General Relativity and Quantum Mechanics Jointly}\\[2pt]
\large Version 2 (November 2026): controlled results and the manifold's field equations}
\author{Rahul Modak$^{1}$ \quad Dr.\ Rahul Walawalkar$^{2}$\\[2pt]
\small $^1$Bayesian Cybersecurity Pvt.\ Ltd., Mumbai/Thane, India \quad
$^2$Carnegie Mellon University; NETRA; Caret Capital\\
\small Code: \url{https://github.com/RahulModak74/QUANTUM_GRAVITY_WITH_MW} (folder \texttt{Nov-2026\_tests})\\
\small Supersedes the March 2026 version of this note}
\date{}

\begin{document}
\twocolumn[
\begin{@twocolumnfalse}
\maketitle
\begin{framed}
\small\textbf{Disclaimer.} We do not claim to have solved quantum gravity. The training data are
synthetic, and the results describe a learned object: a manifold that holds GR and QM content
together. The field equations in Section~5 are the equations \emph{that manifold} satisfies,
written in Einstein's form. They are not offered as laws of nature.
\end{framed}
\begin{abstract}
A Bayesian variational autoencoder (VAE) with a 20-dimensional latent space is trained on
30{,}000 synthetic spacetime states from six regimes (dust, radiation, vacuum, black hole,
wormhole, Hawking radiation). Each row carries GR and QM quantities drawn from one joint sample.
GR and QM constraints enter as priors, and the measurement process enters as a Student-t prior
with $\nu=0.8$. Version 2 replaces the March 2026 results with controlled runs on corrected code
and adds the manifold's field equations. \textbf{(1)} One latent point decodes GR and QM
quantities together: Einstein trace residual 1.3\%, Hawking $T_H\propto M^{-0.93}$ ($R^2=0.83$).
\textbf{(2)} The M-W priors make the joint fit faithful: without them the Hawking fit drops to
$R^2=0.42$. \textbf{(3)} The manifold carries the joint, not the columns: shuffling rows within
each regime raises the Einstein residual to 60\%. \textbf{(4)} The manifold's curvature obeys
$G_{ab}=\kappa T^{\rm MW}_{ab}$ exactly ($\GaussA$), with gravity and quantum sources entering
additively and $\nabla^aT^{\rm MW}_{ab}=0$ ($\BianA$).
\end{abstract}
\vspace{1em}
\end{@twocolumnfalse}
]

\section{Motivation}
GR requires a smooth differentiable manifold; QM requires probabilistic behaviour that is
non-smooth at the Planck scale. Existing approaches add dimensions (string theory), discretise
spacetime (LQG, CDT) or restrict to anti-de Sitter geometry (AdS/CFT). The M-W framework,
developed for physics-informed anomaly detection in battery electrochemistry, takes a different
route: it learns one joint density over both theories' quantities, with their constraints as
Bayesian priors. Coupling theories then becomes density estimation on one manifold.

\section{The Student-t($\nu=0.8$) measurement prior}
\[
p(x\,|\,\nu)=\frac{\Gamma(\frac{\nu+1}{2})}{\sqrt{\nu\pi}\,\Gamma(\frac{\nu}{2})}
\Bigl(1+\frac{x^2}{\nu}\Bigr)^{-\frac{\nu+1}{2}},\quad p\sim|x|^{-(\nu+1)}.
\]
For $\nu>2$ the variance is finite; for $1<\nu\le2$ the mean exists but the variance does not;
for $\nu\le1$ neither exists. At $\nu=0.8$ the density tail decays as $|x|^{-1.8}$: normalisable,
with no finite moments. We use this regime for the Planck/topology flags (prior C6), so rare
discontinuous events are not penalised as outliers. The other five priors (Einstein trace,
wavefunction normalisation, spacetime uncertainty, energy condition, metric signature) are
Gaussian. The choice $\nu=0.8$ is a modelling decision; we do not claim it is derived.

\section{Construction (v5 code)}
Each row is one draw from one joint distribution: the physical parameters are sampled once and
every column is derived from them (the March code drew the black-hole mass twice). The priors read
their columns by name (the March code read three of them one column off). Encoder and decoder have
hidden widths 384/256/192 with LayerNorm and tanh, latent dimension 20, and 150 epochs. A plain VAE
without priors and a row-shuffled control are trained the same way.

\section{Results}
\begin{table}[h]\centering\small
\begin{tabular}{@{}lccc@{}}\toprule
Measure (decoded test states) & M-W & plain & shuffled\\\midrule
Einstein trace residual & 1.3\% & 1.1\% & 59.6\%\\
Hawking slope ($-1$ exact) & $-0.93$ & $-0.66$ & --\\
Hawking $R^2$ & 0.83 & 0.42 & 0.19\\\bottomrule
\end{tabular}
\caption{Results 1--3 (source: \texttt{results.json}).}
\end{table}
Results 1--3 are reported in full in the Nov-2026 paper. The shuffle control keeps every value of
every column but breaks which values occur together. Because it destroys both relations, they
are carried by the joint structure the manifold learns, not by the individual columns.

\section{The manifold's field equations}
With the decoder $f$ as an immersion of the \ActA{} active latent axes into standardised data space,
the pullback metric $g_{ab}=\langle\partial_af,\partial_bf\rangle$, its connection and the second
fundamental form $\mathrm{II}_{ab}$ are computed exactly. The Gauss equation gives
\begin{align*}
G_{ab}&=\kappa T^{\rm MW}_{ab},\\
\kappa T^{\rm MW}_{ab}&=\langle H,\mathrm{II}_{ab}\rangle-g^{cd}\langle\mathrm{II}_{ac},\mathrm{II}_{bd}\rangle\\
&\quad-\tfrac12g_{ab}(|H|^2-|\mathrm{II}|^2),
\end{align*}
and $T^{\rm MW}$ splits exactly into gravity, quantum, horizon, coordinate and Planck-flag terms.
On \NptsA{} unseen test states, the Riemann tensor computed from the connection and the one
computed from $\mathrm{II}$ agree to $\GaussA$; $\nabla^aG_{ab}=0$ holds to $\BianA$. In the
black-hole regime the shares of $|R|$ are gravity $\ShGrBlackholeA\%$, quantum $\ShQuBlackholeA\%$,
horizon $\ShHoBlackholeA\%$; in dust the quantum columns dominate ($\ShQuDustA\%$). A geodesic between
black holes of mass $\GeoMsA$ and $\GeoMeA$ passes through black holes of uniformly increasing
mass. The derivation and all checks are in \emph{Field Equations of the M-W Manifold, Version 2}.

\section{Structural comparison}
\begin{table}[h]\centering\footnotesize\setlength{\tabcolsep}{4pt}
\begin{tabular}{@{}lcccc@{}}\toprule
Approach & Extra dim. & Discretise & AdS only & Single obj.\\\midrule
String theory & Yes & No & Partial & No\\
LQG & No & Yes & No & No\\
CDT & No & Yes & No & No\\
AdS/CFT & Yes & No & Yes & No\\
\textbf{M-W} & \textbf{No} & \textbf{No} & \textbf{No} & \textbf{Yes}\\\bottomrule
\end{tabular}
\end{table}
The other approaches are full theories of quantum gravity; M-W is compared on structure only.
\textbf{M-W removes the coupling step that PDE-based methods require.}

\section{What changed from Version 1}
Version 1 (March 2026) listed five ``emergent'' properties. Version 2 keeps only what the corrected
v5 code and controls support: results 1--3 above, plus the field equations. Properties not
re-tested here (the LIGO waveform comparison, the latent-separation claim as stated, and the
Planck-discreteness interpretation of the tail exponent) are not carried forward. The Version-1
figure and its 3.3\% Einstein residual came from the v4 pipeline with the column-index error.

% Shared section, \input by all three v2 notes. Numbers come from numbers.tex.
\section{Falsifiable predictions}
\label{sec:falsify}
The M-W framework makes a prediction that standard quantum mechanics does not (Appendix~A,
January 2026). Standard QM has $\Delta x\,\Delta p\ge\hbar/2$ and $\Delta E\,\Delta t\ge\hbar/2$, but no
lower bound on joint \emph{position--time} localisation. M-W predicts
\begin{equation}
\Delta x\,\Delta t\;\ge\;\ell_P^{2}\bigl(1+\beta R\,\ell_P^{2}\bigr),\qquad \ell_P=\sqrt{\hbar G/c^{3}},
\label{eq:xt}
\end{equation}
with $\beta$ a dimensionless geometric coupling expected to be of order unity. The prediction can
be falsified at three independent levels, and failure at any one falsifies it:
\begin{enumerate}\itemsep0pt
\item \textbf{Existence:} is there any lower bound on $\Delta x\,\Delta t$? Standard QM predicts none.
\item \textbf{Scale:} if a bound exists, is it set by the Planck area $\ell_P^{2}$?
\item \textbf{Geometry:} does the bound grow with curvature as $1+\beta R\ell_P^{2}$?
\end{enumerate}
The defining experimental signature, playing the role Mercury's perihelion played for GR, would be
a fundamental lower bound on joint spatial--temporal localisation that depends measurably on
spacetime curvature. Candidate tests: joint localisation with ultra-cold atoms or trapped ions plus
precision clocks, compared across gravitational environments. Further channels: corrections to
the Hawking temperature, logarithmic corrections to black-hole entropy, and curvature-dependent
primordial fluctuations.

\paragraph{What the manifold can and cannot say.} On the \UStatesA{} test states the decoded bound
$\Delta x\Delta t\ge\ell_P^{2}$ holds for $\UHeldA$ ($\UBoundA\%$; no-prior model $\UBoundB\%$). This is a
consistency check, not a test, because the bound is built into the training data and into prior C3.
With the exact Ricci scalar, the effective coupling on the manifold is $\beta_{\rm eff}\approx0$
(interquartile range $[\UBetaLoA,\UBetaHiA]$; correlation of the excess with $R$: $\UCorrA$), as it must
be, since the synthetic data contain no curvature dependence. The earlier value
$\beta_{\rm eff}\approx3.24$ was computed with the Version-1 curvature proxy, which is not a curvature
(see \emph{Field Equations, Version 2}), and is withdrawn. $\beta$ is therefore a prediction to be
settled by experiment, not a number the manifold supplies. To test the geometric level on a
manifold, the training data must be generated with a curvature-dependent bound and $\beta$ recovered
blind.

\paragraph{Falsification tests already run on the manifold.} Two controls could have failed and did
not. (i) Shuffling columns within each regime keeps every value but breaks the joint rows; if the
GR+QM relations lived in single columns they would survive, but the Einstein residual rises from
1.3\% to 59.6\% and the Hawking fit collapses ($R^2$ 0.83 $\to$ 0.19). (ii) Removing the M-W priors
lowers the Hawking fit from $R^2=0.83$ to $0.42$. A third check falsifies extraction errors: the
Gauss and Bianchi identities hold to $\GaussA$ and $\BianA$ with exact geometry, and fail for the
Version-1 extraction (its Ricci trace is identically zero).

\paragraph{$\nu>1$ test (run, criteria pre-registered).} The March 2026 prediction was that replacing
the Student-t($\nu=0.8$) prior with $\nu>1$ keeps GR but destroys the Planck/topology structure.
We retrained v5 on identical data for $\nu\in\{0.8,1.5,3.0\}$ with two seeds each, under a decision
rule fixed before training. GR was preserved: Einstein residuals were 0.9--1.3\% for every run, so
the March claim of a 0\% residual for $\nu>1$ is not reproduced. The Planck structure was \emph{not}
degraded: topology-flag accuracy (0.9989--0.9990) and flag error (0.010--0.011) were unchanged, and
only wormhole recall fell slightly (0.976--0.983 against 0.985). Under the pre-registered rule the
prediction is therefore \textbf{falsified on this toy setup}. This does not show that $\nu=0.8$ is
worse: all values perform comparably, and $\nu=0.8$ remains a modelling choice. Outside the
pre-registered rule, the Hawking fit was weaker for $\nu=3.0$ ($R^2$ 0.52, 0.57 against 0.67--0.83
for $\nu=0.8$). This is a hint for future data sets that contain genuine heavy-tailed measurement
events, not a result.


\section{Limitations}
Training data are synthetic and in Planck units over a moderate range. Latent coordinates are not
spacetime coordinates, and the field equations hold for the learned manifold with the Euclidean
metric of standardised data. Results are from one training seed per configuration.

\section*{Reproduce}
\small Results 1--3: \texttt{qg\_v5\_experiments.py} in \texttt{Nov-2026\_tests/code}.
Field equations: \texttt{mw\_exact\_geometry.py} in \texttt{EXPLORATORY\_PHYSICS/code}. Code
implementation was assisted by AI coding tools; modelling choices and claims are the authors'.
\end{document}
